\title{DeepSeekMath: Pushing the Limits of Mathematical Reasoning in Open Language Models}

\begin{abstract}
Mathematical reasoning remains a formidable hurdle for language models because of its inherent complexity and highly structured characteristics. This study presents DeepSeekMath~7B, which extends the pre-training of DeepSeek-Coder-Base-v1.5 7B by utilizing 120B math-focused tokens derived from Common Crawl, alongside data in natural language and code. DeepSeekMath~7B secures a remarkable 51.7\% on the MATH benchmark at a competition level, accomplishing this without the support of external toolkits or ensemble methods, and thus approaches the results of Gemini-Ultra and GPT-4. Employing self-consistency with 64 output samples, DeepSeekMath~7B attains 60.9\% on the MATH benchmark. Two central factors drive DeepSeekMath’s proficiency in mathematical reasoning: Firstly, we capitalize on vast public web resources through a precisely curated data filtering and selection process. Secondly, we propose Group Relative Policy Optimization (GRPO), a novel adaptation of Proximal Policy Optimization (PPO), which improves both mathematical reasoning skills and memory efficiency compared to standard PPO.
\end{abstract}

\section{Introduction}

Large language models (LLMs) have dramatically transformed methods for addressing mathematical reasoning in artificial intelligence, leading to breakthroughs on established benchmarks in quantitative \citep{MATH} and geometric reasoning \citep{trinh2024solving}. In addition, these models have shown considerable value in aiding humans to tackle advanced mathematical challenges \citep{tao}. Nevertheless, the most advanced models, including GPT-4 \citep{gpt4} and Gemini-Ultra \citep{gemini}, remain inaccessible to the public, while existing open-source alternatives exhibit significantly lower levels of performance.

This paper presents DeepSeekMath, a language model specialized for mathematical tasks, which surpasses the mathematical proficiency of current open-source models and nearly matches GPT-4's performance on standard academic evaluations. To develop this capability, we construct the DeepSeekMath Corpus, a comprehensive and high-quality pre-training dataset containing 120B mathematical tokens. The dataset is assembled from the Common Crawl (CC) by leveraging a classifier built on the fastText framework \citep{joulin2016fasttext}. In its first phase, the classifier is trained with examples from OpenWebMath \citep{openwebmath} as positives, while a broad assortment of unrelated web documents functions as negatives. The classifier then identifies additional relevant data within the CC, and this output is subject to further refinement through manual annotation. After augmentation with these curated instances, the classifier undergoes retraining to enhance its accuracy. Our empirical results demonstrate the exceptional quality of this large-scale corpus: DeepSeekMath-Base 7B attains 64.2\% accuracy on GSM8K \citep{gsm8k} and 36.2\% on the MATH competition dataset \citep{MATH}, thus exceeding the performance of Minerva 540B \citep{minerva}. Notably, due to the multilingual nature of the DeepSeekMath Corpus, we observe improved results on Chinese mathematics benchmarks \citep{wei2023cmath,agieval}. We consider our work in curating mathematical data to be a valuable foundation for future advancements, and recognize considerable opportunities for further progress in this area.

DeepSeekMath-Base is built upon DeepSeek-Coder-Base-v1.5 7B \citep{deepseek-coder}, following the observation that models pretrained on code tend to offer a more effective foundation than those starting from a standard large language model. In addition, we find that mathematical pretraining not only enhances the model’s mathematical reasoning, but also yields improvements on general reasoning benchmarks, such as MMLU \citep{mmlu} and BBH \citep{bbh}.

Following the pre-training stage, we fine-tune DeepSeekMath-Base with mathematical instruction datasets incorporating chain-of-thought \citep{cot}, program-of-thought \citep{pot,pal}, and tool-augmented reasoning strategies \citep{tora}. This produces DeepSeekMath-Instruct 7B, which outperforms all other 7B models and even competes with 70B scale open-source instruction-tuned systems.

We further present Group Relative Policy Optimization (GRPO), an alternative reinforcement learning (RL) algorithm related to Proximal Policy Optimization (PPO) \citep{schulman2017proximal}. GRPO dispenses with the use of a critic model; instead, it leverages group-wise baseline estimation from aggregate group scores, thereby dramatically lowering computational costs during training. When applied to a select subset of English instruction tuning data, GRPO demonstrates notable gains over DeepSeekMath-Instruct, with marked performance improvements on both in-domain (GSM8K: 82.9\% $\rightarrow$ 88.2\%, MATH: 46.8\% $\rightarrow$ 51.7\%) and out-of-domain tasks (for example, CMATH: 84.6\% $\rightarrow$ 88.8\%) during RL training. Moreover, we introduce a cohesive conceptual framework to situate methods such as Rejection Sampling Fine-Tuning (RFT) \citep{yuan2023scaling}, Direct Preference Optimization (DPO) \citep{dpo}, PPO, and GRPO within a unified perspective. Analysis under this framework reveals that all these techniques can be interpreted as forms of direct or simplified RL. We perform comprehensive empirical studies—contrasting online versus offline learning, process versus outcome supervision, and single-step versus iterative RL—delving into the critical components underpinning this framework. Lastly, we elucidate the mechanisms by which RL enhances instruction-tuned model performance, and outline prospective avenues for advancing RL effectiveness in light of our unified paradigm.

\subsection{Contributions}
This work introduces advancements in scalable math pre-training, complemented by an in-depth investigation into reinforcement learning techniques.

\textbf{Math Pre-Training at Scale}

\begin{itemize}[topsep=0pt]
    \item We demonstrate that Common Crawl, a publicly available web dataset, harbors substantial mathematical content. By applying a rigorous and systematic data selection process, we develop the DeepSeekMath Corpus—a large-scale, high-quality dataset comprising 120B tokens sourced from mathematically relevant web pages. This dataset surpasses the math web page collections of Minerva \citep{minerva} by nearly sevenfold and is nine times larger than the recent OpenWebMath corpus \citep{openwebmath}.

    \item DeepSeekMath-Base 7B, our pre-trained foundational model, matches the performance of Minerva 540B \citep{minerva}, implying that model size is not the only determinant of mathematical reasoning abilities. Notably, a relatively compact model, when trained on high-caliber data, can achieve impressive results.

    \item We report outcomes from our math pre-training studies. We observe that performing code pre-training before math pre-training enhances the model’s effectiveness at solving mathematical problems, irrespective of tool use. This partially addresses the frequently debated issue: \textit{does code pre-training strengthen reasoning skills?} Our results suggest it does, particularly in mathematical contexts.

    \item Despite the prevalent practice of incorporating arXiv papers in math pre-training pipelines, our experiments show no significant gains on any of the mathematical benchmarks employed in this study.
\end{itemize}

\textbf{Exploration and Analysis of Reinforcement Learning}

\begin{itemize}[topsep=0pt]
    \item We propose Group Relative Policy Optimization (GRPO), a reinforcement learning method that operates without a critic model by leveraging group scores to estimate the baseline, thereby achieving substantial reductions in computational overhead compared to Proximal Policy Optimization (PPO).
    \item Our results show that GRPO considerably boosts the capabilities of our instruction-tuned system, DeepSeekMath-Instruct, even when only instruction-tuning data is used. Additionally, we note that reinforcement learning with GRPO yields improved results on tasks outside the training distribution.
    \item We offer a cohesive analytical framework encompassing approaches such as RFT, DPO, PPO, and GRPO. Moreover, we present thorough empirical analyses, such as comparisons between online and offline training, outcome- versus process-based supervision, as well as single-turn and iterative reinforcement learning, in order to systematically elucidate key components of this framework.
    \item Building upon our integrative framework, we investigate why reinforcement learning is effective for these tasks and outline several promising research avenues aimed at further enhancing LLM reinforcement learning.
\end{itemize}

\subsection{Summary of Evaluations and Metrics}

\begin{itemize}[topsep=0pt]
    \item \textbf{English and Chinese Mathematical Reasoning}: We perform in-depth evaluations of our models using both English and Chinese test sets that encompass mathematics questions at educational levels ranging from primary school to university. For English tasks, benchmarks considered are GSM8K \citep{gsm8k}, MATH \citep{MATH}, SAT \citep{llemma}, OCW Courses \citep{minerva}, and MMLU-STEM \citep{mmlu}. For Chinese, we assess using MGSM-zh \citep{mgsm}, CMATH \citep{wei2023cmath}, Gaokao-MathCloze \citep{agieval}, and Gaokao-MathQA \citep{agieval}. Our evaluation measures both autonomous problem-solving with written solutions—without recourse to tools—and problem-solving with Python code.

    DeepSeekMath-Base exhibits performance on par with the proprietary Minerva 540B \citep{minerva} on English datasets and consistently outperforms all open-source foundation models, including Mistral 7B \citep{mistral} and Llemma-34B \citep{llemma}, independent of their exposure to mathematical pre-training and frequently by a large extent. Its strength on Chinese mathematical tasks stands out as well, potentially due to our data collection strategy that augments English resources with extensive high-quality non-English content, diverging from prior practices \citep{minerva,llemma}. Through the integration of instruction-tuning and reinforcement learning, both DeepSeekMath-Instruct and DeepSeekMath-RL attain state-of-the-art accuracy, achieving over 50\% on the challenging, competition-grade MATH dataset—a milestone among open-source projects.
\end{itemize}

\item \textbf{Formal Mathematics}: DeepSeekMath-Base is assessed on its informal-to-formal theorem proving abilities using the task established in \citep{dsp_proof}, applied to miniF2F \citep{minif2f} with Isabelle \citep{isabelle} selected as the proof assistant. Results indicate that DeepSeekMath-Base achieves robust autoformalization performance in few-shot settings.
\item \textbf{Natural Language Understanding, Reasoning, and Code}: To thoroughly assess the models' capabilities in understanding, reasoning, and coding, we evaluate DeepSeekMath-Base on the Massive Multitask Language Understanding (MMLU) benchmark \citep{mmlu}, which features 57 multiple-choice tasks from a variety of disciplines; BIG-Bench Hard (BBH) \citep{bbh}, consisting of 23 tasks primarily demanding complex multi-step reasoning; as well as HumanEval \citep{codex} and MBPP \citep{mbpp}, both recognized for measuring code generation proficiency. Pre-training with mathematical data is found to enhance both language understanding and reasoning effectiveness.

\section{Math Pre-Training}

\subsection{Data Collection and Decontamination} This section details the construction of the DeepSeekMath Corpus using data sourced from Common Crawl. Figure \ref{fig:data_collect} illustrates the iterative pipeline adopted to systematically amass a large-scale mathematical dataset, beginning with a compact, high-quality math-oriented seed set. Notably, this methodology is flexible and could be extended to other specialized domains, such as programming.

Initially, OpenWebMath \citep{openwebmath}—a curated selection of quality mathematical web resources—is used as the starting seed corpus. A fastText model \citep{joulin2016fasttext} is then trained to retrieve additional web pages similar in content to OpenWebMath. To accomplish this, 500,000 samples are drawn at random from the seed set as positive examples, while another 500,000 web pages from Common Crawl serve as negative examples. The training employs an open-source fastText implementation\footnote{\url{https://fasttext.cc}}, with parameters set to a vector size of 256, learning rate at 0.1, word n-grams up to length 3, a minimum word frequency threshold of 3, and 3 epochs of training. The original Common Crawl dataset is pruned by applying both URL-based deduplication and near-deduplication, which reduces the pool to 40B HTML pages. Mathematical content is subsequently extracted from this deduplicated dataset using the trained fastText model. Pages are then ranked by the scores output by fastText to filter out lower-quality entries, retaining only the highest scoring content. Pre-training experiments on datasets containing the top 40B, 80B, 120B, and 160B tokens are used to determine the optimal volume to preserve. For the first iteration, the top 40B tokens are retained.

Following the initial phase of data acquisition, a substantial portion of mathematical web pages remains undiscovered. This limitation stems chiefly from the inadequate diversity of positive examples within the initial fastText training dataset. To address this, we incorporate additional mathematical web resources to further augment the seed corpus, thereby facilitating enhanced optimization of the fastText model. Our procedure begins by partitioning the entire Common Crawl dataset into non-overlapping domains, where a domain encompasses all web pages sharing the same base URL. For each domain, we assess the proportion of web pages captured during the first iteration. Domains in which more than 10\% of pages have been collected are flagged as mathematics-oriented (e.g., \url{mathoverflow.net}).

We then undertake a manual review to annotate URLs in these domains that correspond to mathematical content (for example, \url{mathoverflow.net/questions}). Pages associated with these mathematical URLs that were not previously collected are added to the seed corpus. By this method, we acquire an expanded set of positive instances, enabling subsequent training of a more robust fastText model that can recall a broader range of mathematical pages in future iterations. After four collection cycles, our process results in 35.5M mathematical web pages with a cumulative total of 120B tokens. By the fourth iteration, we observe that almost 98\% of the corpus has already been harvested in the prior iteration, prompting us to conclude the collection process.

To minimize the risk of contaminating evaluation benchmarks, we employ the filtering strategy used in \cite{deepseek-coder}, removing web pages that contain content from prominent mathematical benchmarks in English, such as GSM8K \citep{gsm8k} and MATH \citep{MATH}, as well as Chinese datasets like CMATH \citep{wei2023cmath} and AGIEval \citep{agieval}. The filtration method is as follows: any segment of text that includes a 10-gram sequence exactly matching a substring from the benchmarks is excluded from the training dataset. For benchmark excerpts shorter than 10 grams but at least 3 grams, we utilize exact substring matching to filter out affected pages.

\subsection{Quality Assessment of the DeepSeekMath~Corpus}

To assess the relative quality of the DeepSeekMath~Corpus, we conduct pre-training studies benchmarking it against several newly available math-focused corpora:

\begin{itemize}[topsep=0pt]
    \item \textbf{MathPile} \citep{mathpile}: an 8.9B-token, multi-source corpus drawing content from textbooks, Wikipedia, ProofWiki, CommonCrawl, StackExchange, and arXiv, with arXiv constituting more than 85\% of the data;
    \item \textbf{OpenWebMath} \citep{openwebmath}: a 13.6B-token dataset curated from CommonCrawl via mathematical content filtering;
    \item \textbf{Proof-Pile-2} \citep{llemma}: a composite mathematical dataset formed by aggregating OpenWebMath, AlgebraicStack (containing 10.3B tokens of mathematical code), and 28.0B tokens of arXiv publications. In our experiments on Proof-Pile-2, we follow the arXiv:Web:Code ratio of 2:4:1 as proposed by \cite{llemma}.
\end{itemize}

\subsubsection{Training Setting}
\label{sec:quality-policy}
We perform mathematical domain adaptation on a general-purpose language model consisting of 1.3B parameters, following the same architecture as the DeepSeek LLMs \citep{deepseek-llm}, henceforth referred to as DeepSeek-LLM 1.3B. Separate instances of the model are further trained on each specific mathematics dataset, each for a total of 150B tokens. Training is implemented via the lightweight and efficient HAI-LLM \citep{haillm} framework. In alignment with the training methodologies adopted in DeepSeek LLMs, we utilize the AdamW optimizer \citep{adamW}, setting $\beta_1=0.9$, $\beta_2=0.95$, and a $\mathrm{weight\_decay}=0.1$. A multi-stage learning rate schedule is employed: the learning rate ramps up to its maximum after a warmup phase of 2,000 steps, reduces to 31.6\% of the peak at the 80\% completion mark, and then further declines to 10.0\% of the peak by 90\% of the training process. The peak learning rate is configured at 5.3e-4, with a total batch size of 4M tokens and a sequence length of 4K tokens.

\subsubsection{Evaluation Results}

\textbf{The DeepSeekMath~Corpus offers high-quality, large-scale, and multilingual mathematical data, surpassing existing resources in these aspects.}

\begin{itemize}[topsep=0pt]
    \item \textbf{High-quality}:
    We assess downstream capabilities using 8 mathematical reasoning benchmarks under few-shot chain-of-thought prompting \cite{cot}. According to Table \ref{tab:corpora_comparison}, the model trained with the DeepSeekMath~Corpus consistently surpasses others in performance. As illustrated in Figure \ref{fig:corpora_comparisons}, models trained on DeepSeekMath significantly outperform those trained on Proof-Pile-2 after processing 50B tokens (equivalent to a full epoch for Proof-Pile-2), highlighting the superior average quality of the DeepSeekMath data.
    \item \textbf{Multilingual}:
    The DeepSeekMath~Corpus includes content in several languages, with English and Chinese comprising the majority. Table \ref{tab:corpora_comparison} reveals that leveraging the DeepSeekMath~Corpus during training produces notable gains in mathematical reasoning for both English and Chinese tasks. By contrast, current mathematical datasets, predominantly in English, display restricted enhancements or even adverse effects on Chinese mathematical reasoning tasks.
    \item \textbf{Large-scale}:
    DeepSeekMath~Corpus is multiple times larger than comparable mathematical corpora available. Figure \ref{fig:corpora_comparisons} shows that training DeepSeek-LLM 1.3B with this extensive corpus results in more pronounced and sustained improvements, reflected in a steeper learning trajectory. The baseline corpora, being much smaller, have been cycled through numerous times, leading to an early plateau in their model performance.
\end{itemize}

\subsection{Training and Evaluating DeepSeekMath-Base 7B}

This section details the development of DeepSeekMath-Base 7B, a foundational model with advanced reasoning capabilities tailored to mathematics. DeepSeekMath-Base 7B is initialized from DeepSeek-Coder-Base-v1.5 7B \citep{deepseek-coder} and undergoes training on a dataset comprising 500B tokens. The training corpus is distributed as follows: 56\% sourced from the DeepSeekMath Corpus, 4\% from AlgebraicStack, 10\% from arXiv, 20\% consisting of Github code, while the remaining 10\% originates from Common Crawl in English and Chinese natural language. The majority of training procedures follow the configurations outlined in Section \ref{sec:quality-policy}, with two key adjustments: the learning rate peaks at 4.2e-4 and a batch size of 10M tokens is used.

A thorough evaluation of DeepSeekMath-Base 7B is performed, emphasizing its proficiency in autonomously generating complete mathematical solutions, leveraging tools to resolve mathematical problems, and performing formal theorem proofs. In addition to its mathematical competence, the model’s general capabilities are also assessed, encompassing tasks in natural language understanding, reasoning, and programming.

\paragraph{Mathematical Problem Solving with Step-by-Step Reasoning}

To measure DeepSeekMath-Base's aptitude in mathematical problem solving, we employ few-shot chain-of-thought prompting \citep{cot} on eight distinct benchmarks spanning both English and Chinese. These benchmarks evaluate abilities in quantitative reasoning (such as GSM8K \citep{gsm8k}, MATH \citep{MATH}, and CMATH \citep{wei2023cmath}) and multiple-choice formats (such as MMLU-STEM \citep{mmlu} and Gaokao-MathQA \citep{agieval}), encompassing a broad spectrum of mathematical domains from primary school concepts to undergraduate-level topics.

As detailed in Table \ref{tab:base_math_cot}, DeepSeekMath-Base 7B attains top results across all eight benchmarks, surpassing other open-source base models, including the popular Mistral 7B \citep{mistral} and the more recently introduced Llemma 34B \citep{llemma}, which benefited from mathematical pretraining on Proof-Pile-2 \citep{llemma}. Particularly noteworthy is its performance on the competitive MATH dataset, where DeepSeekMath-Base exceeds the best open-source base model by a margin greater than 10\% in absolute terms. Furthermore, it outperforms Minerva 540B \citep{minerva}—a much larger closed-source base model (77 times the size) that is built upon PaLM \citep{palm} and further adapted with mathematical texts.

\paragraph{Mathematical Problem Solving with Tool Use}  
For program-assisted mathematical problem solving, we assess models on GSM8K and MATH using few-shot program-of-thought prompting strategies \citep{pot,pal}. Each model is tasked to resolve problems by generating Python scripts, utilizing libraries such as \textit{math} and \textit{sympy} as needed for complex calculations, and answers are determined by executing the program output. According to Table \ref{tab:base_math_pot_proof}, DeepSeekMath-Base 7B outstrips the former leading Llemma 34B on these tasks.

\paragraph{Formal Mathematics}

Automation in formal proof construction advances both the reliability and efficiency of verifying mathematical results, a topic of increasing focus in recent studies. We measure DeepSeekMath-Base 7B's capacity in informal-to-formal proof generation, as in \citep{dsp_proof}, wherein the model produces formal proofs given an informal problem statement, a corresponding formal expression, and an informal argument. The evaluation is carried out on miniF2F \citep{minif2f}, which targets formalizations at Olympiad difficulty, with few-shot prompting to construct Isabelle proofs per instance. Consistent with the approach of \cite{dsp_proof}, models are used to draft proof outlines, while Sledgehammer \citep{sledgehammer}, an automated proof assistant, is employed to complete gaps. Results in Table \ref{tab:base_math_pot_proof} show DeepSeekMath-Base 7B to be highly effective at proof formalization.

\paragraph{Natural Language Understanding, Reasoning, and Code}

Model capability in language understanding is tested on MMLU \citep{mmlu}, for reasoning on BBH \citep{bbh}, and coding performance is evaluated with HumanEval \citep{codex} and MBPP \citep{mbpp}. Table \ref{tab:understanding_reasoning_code} indicates that DeepSeekMath-Base 7B achieves considerable gains on MMLU and BBH relative to its predecessor DeepSeek-Coder-Base-v1.5 \citep{deepseek-coder}, confirming that math-specific training enriches both comprehension and reasoning. Notably, with continual code token exposure, DeepSeekMath-Base 7B preserves the coding capabilities of DeepSeek-Coder-Base-v1.5 on both coding benchmarks. Taken together, DeepSeekMath-Base 7B substantially exceeds the general-purpose Mistral 7B \citep{mistral} on all three reasoning and coding evaluations.

\section{Supervised Fine-Tuning}

\subsection{SFT Data Curation}

We assemble a dataset for mathematical instruction tuning that encompasses both English and Chinese questions across a range of mathematical domains and difficulties. Each question is matched with detailed solutions provided in chain-of-thought (CoT) \citep{cot}, program-of-thought (PoT) \citep{pot,pal}, or tool-integrated reasoning formats \citep{tora}. The curated dataset contains a total of 776K training samples.

\begin{itemize}[topsep=0pt]
    \item \textbf{English mathematical datasets}: We provide tool-integrated solutions to GSM8K and MATH problems, and incorporate selected items from MathInstruct \citep{MathInstruct} as well as training data from Lila-OOD \citep{lila} where either CoT or PoT approaches are employed. The English portion addresses a broad spectrum of mathematical areas, including but not limited to algebra, probability, number theory, calculus, and geometry.
    \item \textbf{Chinese mathematical datasets}: We gather K-12 level Chinese mathematical questions from 76 sub-categories, such as linear equations. Each question is annotated with both CoT-style and tool-integrated reasoning solutions.
\end{itemize}

\subsection{Training and Evaluating DeepSeekMath-Instruct 7B}

Here, we present DeepSeekMath-Instruct 7B, which is adapted for mathematical instruction following the DeepSeekMath-Base architecture. During training, instances are concatenated at random until the input length reaches 4K tokens. The model is optimized over 500 steps, using a batch size of 256, and maintaining a fixed learning rate of 5e-5.

We assess model performance with and without external tool support, using four benchmarks that test quantitative reasoning in both English and Chinese. Comparative analysis includes state-of-the-art models available at the time:

\begin{itemize}[topsep=0pt]
    \item \textbf{Closed-source models} considered are: (1) the GPT series, notably GPT-4 \citep{gpt4} and GPT-4 Code Interpreter~\footnote{\url{https://openai.com/blog/chatgpt-plugins\#\#code-interpreter}}, (2) Gemini Ultra and Pro \citep{gemini}, (3) Inflection-2 \citep{inflection-2}, (4) Grok-1~\footnote{\url{https://x.ai/model-card}}, as well as recently introduced offerings from Chinese firms such as (5) Baichuan-3~\footnote{\url{https://www.baichuan-ai.com}}, (6) and the newest GLM-4~\footnote{\url{https://open.bigmodel.cn/dev/api\#glm-4}}
    in the GLM family \citep{glm}.
\end{itemize}

These models are designed for general applications, with the majority having been subjected to alignment procedures.

\item \textbf{Open-source models} consist of: general-purpose systems such as (1) DeepSeek-LLM-Chat 67B \citep{deepseek-llm}, (2) Qwen 72B \citep{qwen}, (3) SeaLLM-v2 7B \citep{seallm}, and (4) ChatGLM3 6B \citep{chatglm3}. Additionally, some models provide advanced mathematical reasoning capabilities, including: (5) InternLM2-Math 20B~\footnote{\url{https://github.com/InternLM/InternLM-Math}}, which extends InternLM2 via math-specific training and subsequent instruction tuning; (6) Math-Shepherd-Mistral 7B, trained on Mistral 7B \citep{mistral} through PPO with a process-supervised reward signal \citep{schulman2017proximal}; (7) the WizardMath collection \citep{wizardmath}, which bolsters math reasoning in Mistral 7B and Llama-2 70B \citep{llama2} using evolve-instruct (a variation of instruction tuning that leverages AI-generated instructions) and PPO, relying heavily on training examples from GSM8K and MATH; (8) MetaMath 70B \citep{metamath}, derived from Llama-2 70B and further refined with an expanded GSM8K and MATH dataset; (9) ToRA 34B \cite{tora}, which adapts CodeLlama 34B through fine-tuning for mathematical reasoning with tool integration; (10) MAmmoTH 70B \citep{MathInstruct}, consisting of Llama-2 70B subject to instruction tuning via MathInstruct.

Table \ref{tab:sft_rl_math} presents results under a tool-restricted evaluation framework, where DeepSeekMath-Instruct 7B demonstrates robust abilities in stepwise mathematical reasoning. In particular, on the MATH dataset reflecting competition-level problems, our model exceeds the performance of all other open-source models and outperforms most proprietary offerings (e.g., Inflection-2 and Gemini Pro) by at least 9\% in absolute terms. This advantage remains even compared with models with substantially larger capacities (e.g., Qwen 72B) or those extensively optimized for mathematics via reinforcement learning methods (e.g., WizardMath-v1.1 7B). While DeepSeekMath-Instruct achieves results on par with Chinese proprietary systems such as GLM-4 and Baichuan-3 on MATH, it still falls short of matching the achievements of GPT-4 and Gemini Ultra.

When evaluating with tool-assisted natural language and code-based problem-solving permitted, DeepSeekMath-Instruct 7B attains close to 60\% accuracy on MATH—higher than any other open-source alternative. On remaining benchmarks, the model achieves parity with DeepSeek-LLM-Chat 67B, the previous state-of-the-art among open models despite its significantly larger scale.

\section{Reinforcement Learning}

\subsection{Group Relative Policy Optimization} Reinforcement learning (RL) has been established as an effective method to further enhance mathematical reasoning in LLMs after supervised fine-tuning (SFT) \citep{wang2023math,wizardmath}. Here, we detail our new RL approach, Group Relative Policy Optimization (GRPO), which offers both improved efficiency and effectiveness.

\subsubsection{From PPO to GRPO} Proximal Policy Optimization (PPO) \citep{schulman2017proximal}, an actor-critic based RL framework, remains prevalent for RL-based fine-tuning of LLMs \citep{ouyang2022training}. It updates model parameters by maximizing the following objective function:

\begin{equation}
\footnotesize
    \mathcal{J}_{PPO}(\theta) = \mathbb{E}{[q \sim P(Q), o \sim \pi_{\theta_{old}}(O|q)]} \frac{1}{|o|} \sum_{t=1}^{|o|} \min \left[ \frac{\pi_\theta(o_{t} | q, o_{<t})}{\pi_{\theta_{old}}(o_{t} | q, o_{<t})} A_{t}, \text{clip} \left( \frac{\pi_\theta(o_{t} | q, o_{<t})}{\pi_{\theta_{old}}(o_{t} | q, o_{<t})}, 1 - \varepsilon,

In this context, $\pi_{\theta}$ and $\pi_{\theta_{old}}$ denote the updated and prior policy models, respectively. The variables $q$ and $o$ represent questions and their corresponding outputs, with sampling drawn from the question dataset and from the output distribution of the previous policy $\pi_{\theta_{old}}$. The parameter $\varepsilon$ serves as a clipping hyperparameter incorporated into PPO, designed to enhance training stability. The term $A_t$ stands for the advantage estimate, calculated via Generalized Advantage Estimation (GAE) \citep{gae}, utilizing future rewards $\{r_{\ge t}\}$ and the estimated value function $V_{\psi}$. Consequently, PPO necessitates training a value function model jointly with the policy model. To mitigate potential overfitting to the reward model, it is common practice to introduce a per-token KL divergence penalty from a reference policy within the reward at every token \citep{ouyang2022training}, specified as follows:

\begin{equation}
    r_{t} = r_\varphi(q, o_{\le t}) - \beta \log\frac{\pi_{\theta}(o_{t}|q, o_{<t})}{\pi_{ref}(o_{t}|q, o_{<t})},
\label{eq:PPO-reward}
\end{equation}

Here, $r_\varphi$ designates the reward model, $\pi_{ref}$ refers to the reference policy—commonly initialized as the supervised fine-tuned (SFT) model—and $\beta$ is a hyperparameter controlling the KL regularization strength.

Given that PPO's value function is generally an additional network on par with the policy model in terms of capacity, this results in notable computational and memory overheads. Moreover, during reinforcement learning, the value function is used as a baseline for the advantage to reduce variance. In large language model (LLM) training, however, the reward model typically assigns a scalar reward only to the last token, which poses challenges for training a value function that is accurate at all token positions. To address these issues, we introduce Group Relative Policy Optimization (GRPO), illustrated in Figure \ref{fig:grpo}, which sidesteps the need for separate value function learning. Instead, GRPO leverages the mean reward over a set of outputs sampled for the same question as a baseline. Specifically, for every question $q$, a batch of outputs $\{o_1, o_2, \cdots, o_G\}$ is generated from the prior policy $\pi_{\theta_{old}}$, and the policy update objective is formulated as:

\begin{equation}
\footnotesize
\begin{split}
    \mathcal{J}_{GRPO}(\theta) &= \mathbb{E}{[q \sim P(Q), \{o_i\}_{i=1}^G \sim \pi_{\theta_{old}}(O|q)]}  \\
    & \frac{1}{G}\sum_{i=1}^G\frac{1}{|o_i|} \sum_{t=1}^{|o_i|} \left\{ \min \left[ \frac{\pi_\theta(o_{i,t} | q, o_{i,<t})}{\pi_{\theta_{old}}(o_{i,t} | q, o_{i,<t})} \hat{A}_{i,t}, \text{clip} \left( \frac{\pi_\theta(o_{i,t} | q, o_{i,<t})}{\pi_{\theta_{old}}(o_{i,t} | q, o_{i,<t})}, 1 - \varepsilon, 1 + \varepsilon \right)  \hat{A}_{i,t} \right] - \beta \mathbb{D}_{KL}\left[\pi_{\theta} || \pi_{ref}\right]\right\} ,
\end{split}
\label{eq:GRPO-obj}
\end{equation}

The hyperparameters $\varepsilon$ and $\beta$ regulate clipping and KL regularization, respectively. The advantage $\hat{A}_{i,t}$ is defined solely by relative rewards computed within each sample group—a process elaborated in the upcoming subsections. The use of intra-group relative rewards for advantage computation complements the comparative structure of reward models, as these are typically trained on datasets built from pairwise comparisons across answers to the same query. Notably, instead of incorporating the KL regularization into the reward as done in PPO, GRPO introduces KL divergence directly in the loss between the updated policy and the reference policy, which simplifies the calculation of $\hat{A}_{i,t}$. Furthermore, unlike the KL penalty term in (\ref{eq:PPO

\begin{algorithm}[t]
  \small
  \caption{Iterative Group Relative Policy Optimization}
  \textbf{Input} initial policy model $\pi_{\theta_{\text{init}}}$; reward models $r_\varphi$; task prompts $\mathcal{D}$; 
  hyperparameters $\varepsilon$, $\beta$, $\mu$
  \begin{algorithmic}[1]
    \State Initialize the policy model: $\pi_\theta \leftarrow \pi_{\theta_{\text{init}}}$
    \For{iteration = 1, \dots, I}
       \State Set the reference model: $\pi_{ref} \leftarrow \pi_{\theta}$
      \For{step = 1, \dots, M}
      \State Draw a batch $\mathcal{D}_b$ from $\mathcal{D}$
      \State Store current policy as $\pi_{\theta_{old}} \leftarrow \pi_{\theta}$ 
      \State For every question $q \in \mathcal{D}_b$, generate $G$ responses $\{o_i\}_{i=1}^G$ from $\pi_{\theta_{old}} (\cdot \mid q)$
      \State For each generated output $o_i$, evaluate reward $r_i$ by applying $r_{\varphi}$
      \State Employ group-based relative advantage calculation to estimate $\hat{A}_{i,t}$ for each token $t$ in $o_i$
      \For{GRPO iteration = 1, \dots, $\mu$}
        \State Optimize $\pi_{\theta}$ by maximizing the GRPO loss as described in Equation \ref{eq:GC-GRPO}
      \EndFor
    \EndFor \State Refine $r_\varphi$ with continual learning using a replay buffer. 
    \EndFor 
  \end{algorithmic}
  \textbf{Output} $\pi_\theta$
  \label{alg:iter-grpo}
\end{algorithm}

\subsubsection{Outcome Supervision RL with GRPO} More precisely, for a given question $q$, we sample a collection of responses $\{o_1, o_2, \cdots, o_G\}$ from the frozen policy $\pi_{\theta_{old}}$. Each response is scored with the reward model, resulting in $G$ reward values $\mathbf{r}=\{r_1, r_2, \cdots, r_G\}$. The reward set is then standardized by subtracting the mean and dividing by the standard deviation within the group. With outcome supervision, the normalized terminal reward for each response $o_i$ is assigned as the advantage $\hat{A}_{i, t}$ to all tokens $t$ in the sequence: $\hat{A}_{i, t} = \widetilde{r}_i = \frac{r_i- {\rm mean}(\mathbf{r})}{{\rm std}(\mathbf{r})}$. Policy updates are conducted to maximize the objective in equation (\ref{eq:GRPO-obj}).

\subsubsection{Process Supervision RL with GRPO} Since outcome-based rewards are only attributed at sequence termination, they can lack informativeness or efficacy for complex reasoning tasks. Building on \cite{wang2023math}, we adopt process supervision, where rewards are attributed to each step of the reasoning sequence. Concretely, given $q$ and a group of $G$ sampled outputs $\{o_1, o_2, \cdots, o_G\}$, the process reward model computes a reward for every reasoning step, yielding $\mathbf{R} = \{ \{r_1^{index(1)},\cdots,r_1^{index(K_1)}\}, \cdots,  \{r_G^{index(1)},\cdots,r_G^{index(K_G)}\} \}$, with $index(j)$ denoting the final token index for the $j$-th step, and $K_i$ representing the number of steps in output $o_i$. These stepwise rewards are similarly standardized: $\widetilde{r}_i^{index(j)} = \frac{r_i^{index(j)} - {\rm mean(\mathbf{R})}}{{\rm std(\mathbf{R})}}$.
Advantage values for each token are then computed by summing normalized future step rewards: $\hat{A}_{i, t} = \sum_{index(j) \ge t} \widetilde{r}_i^{index(j)}$,
after which policy improvement follows via maximizing the objective in equation (\ref{eq:GRPO-obj}).

\subsubsection{Iterative RL with GRPO} As reinforcement learning training advances, the original reward model may become inadequate for supervising the continually evolving policy model. Consequently, we investigate an iterative approach to RL using GRPO. In Algorithm \ref{alg:iter-grpo}, this iterative GRPO method involves constructing fresh reward model training sets by sampling outputs from the updated policy model. The existing reward model is refined via a replay strategy that mixes in 10\% of previously seen data. Subsequently, the current policy model serves as the reference model, and it undergoes ongoing training utilizing the latest version of the reward model.

\subsection{Training and Evaluating DeepSeekMath-RL}

Our RL experiments are carried out using DeepSeekMath-Instruct 7B as the base model. The RL training set comprises chain-of-thought formatted queries sampled from the GSM8K and MATH datasets within the SFT data, totaling approximately 144K items. All other SFT-derived questions are omitted to assess how RL influences benchmarks that lack exposure during the RL process. Construction of the reward model's training dataset follows the procedure described in \citep{wang2023math}. The initial reward model is fine-tuned from DeepSeekMath-Base 7B using a 2e-5 learning rate. When applying GRPO, we fix the policy model’s learning rate at 1e-6, and set the KL penalty to 0.04. Each question receives $64$ sampled responses. The output sequence is restricted to a maximum of 1024 tokens, and the training employs a batch size of 1024. Policy model updates occur just once after each round of sampling. We assess DeepSeekMath-RL 7B using the same benchmarks as DeepSeekMath-Instruct 7B. Within this evaluation, tasks from GSM8K and MATH that involve chain-of-thought reasoning are categorized as in-domain, while the remainder of the benchmarks are treated as out-of-domain.

Table \ref{tab:sft_rl_math} presents the comparative results of open- and closed-source models, evaluating their capabilities in both chain-of-thought and tool-integrated reasoning on English and Chinese datasets. Key observations include: 1) DeepSeekMath-RL 7B achieves accuracy rates of 88.2\% on GSM8K and 51.7\% on MATH with chain-of-thought reasoning, outperforming every open-source model from 7B to 70B parameters as well as most closed-source alternatives. 2) Importantly, DeepSeekMath-RL 7B is exclusively trained with chain-of-thought-formatted instruction tuning data from GSM8K and MATH, with DeepSeekMath-Instruct 7B as the initialization point. Even though its training data is narrowly focused, DeepSeekMath-RL 7B consistently surpasses DeepSeekMath-Instruct 7B on all evaluation benchmarks, illustrating the substantial advantages conferred by reinforcement learning.

\section{Discussion}
This section elaborates on insights derived from our pre-training and RL experiments.

\subsection{Lessons Learnt in Pre-Training}

We begin by outlining observations from the pre-training phase. Unless otherwise noted, training adheres to the settings described in Section \ref{sec:quality-policy}. Throughout this section, references to the DeepSeekMath Corpus pertain to an 89B-token dataset compiled during the second round of data collection.

\subsubsection{Code Training Benefits Mathematical Reasoning}

A widely held, though previously unconfirmed, hypothesis posits that training with code enhances a model’s reasoning abilities. Our investigation sheds partial light on this, focusing on the mathematical context: incorporating code in training enhances a model’s proficiency in mathematical reasoning, regardless of tool usage.

To examine the effect of code-based training on mathematical reasoning, we compared two training paradigms: two-stage training and one-stage training, as defined below:

\textbf{Two-Stage Training}

\begin{itemize}[topsep=0pt]
    \item \textbf{Code Training for 400B Tokens $\rightarrow$ Math Training for 150B Tokens}: DeepSeek-LLM 1.3B undergoes pretraining on 400B code tokens, after which it is further trained on 150B tokens of mathematical content;
    \item \textbf{General Training for 400B Tokens $\rightarrow$ Math Training for 150B Tokens}: For comparison, we substitute the initial code token training with 400B general domain tokens (sourced from DeepSeek-AI’s comprehensive general corpus), preceding 150B math tokens, to assess whether code data confers distinct benefits over general data for developing mathematical reasoning capabilities.
\end{itemize}

\textbf{One-Stage Training}

\begin{itemize}[topsep=0pt]
    \item \textbf{Math Training for 150B Tokens}: Here, DeepSeek-LLM 1.3B is exclusively trained on 150B tokens of mathematical data;
    \item \textbf{Training on a mixture of 400B Code Tokens and 150B Math Tokens}: Sequential math training after code pretraining can impair performance on coding. We therefore examine a single-stage approach where code and math tokens are interleaved, exploring whether this joint training not only benefits mathematical reasoning but also curtails the loss in coding ability typically caused by sequential fine-tuning.
\end{itemize}

\paragraph{Results}
The results summarised in Table \ref{tab:code-math} and Table \ref{tab:code-math-general-eval} capture model performance under the above training regimens.

Incorporating code data during training is advantageous for tasks requiring program-based mathematical reasoning, in both sequential and mixed training frameworks. As observed in Table \ref{tab:code-math}, code pretraining alone greatly enhances model performance on GSM8K and MATH tasks when programmatic solutions in Python are required, and subsequent math-specific training brings additional gains. Notably, mixing code and math data in a unified training stage effectively addresses the catastrophic forgetting commonly seen in multi-stage training pipelines, and yields complementary improvements in both code generation (Table \ref{tab:code-math-general-eval}) and program-oriented mathematical reasoning (Table \ref{tab:code-math}).

Code pretraining also yields gains for pure mathematical reasoning (i.e., tasks without external tool usage). In the sequential training scenario, the initial exposure to code prompts moderate improvement, and it further enhances the efficacy of later math-focused training, culminating in optimal results. However, integrating code and math data simultaneously within a single training phase appears to hamper non-tool mathematical reasoning. A likely explanation is that DeepSeek-LLM 1.3B’s limited parameter count inhibits its ability to comprehensively assimilate both domains when learned concurrently.

\subsubsection{ArXiv Papers Appear Ineffectual for Advancing Mathematical Reasoning}

It is standard practice to include arXiv papers as part of the pre-training datasets for mathematical language models \citep{minerva,gpt-f,llemma,mathpile}. Nevertheless, comprehensive evaluations of their actual contribution to mathematical reasoning abilities remain sparse. Contrary to expectations, our empirical findings indicate that arXiv papers do not significantly enhance mathematical reasoning. We conducted experiments utilizing various model scales—specifically, DeepSeek-LLM 1.3B and DeepSeek-Coder-Base-v1.5 7B \citep{deepseek-coder}—with arXiv datasets that underwent distinct preprocessing techniques:

\begin{itemize}[topsep=0pt]
    \item \textbf{MathPile} \citep{mathpile}: an 8.9B-token dataset curated through heuristics involving cleaning and filtering, where over 85\% originates from scientific arXiv manuscripts;
    \item \textbf{ArXiv-RedPajama} \citep{redpajama}: a 28.0B-token dataset comprised of arXiv LaTeX source files with extraneous elements such as preambles, comments, macros, and bibliographies excluded.
\end{itemize}

For our evaluation, DeepSeek-LLM 1.3B was trained for 150B tokens, and DeepSeek-Coder-Base-v1.5 7B for 40B tokens, utilizing each of the aforementioned arXiv datasets. Across these setups, no marked improvements—or at times, regressions—were observed on a suite of mathematical reasoning benchmarks after exclusive training on arXiv data. These assessments span quantitative reasoning benchmarks including GSM8K and MATH (see Table \ref{tab:arxiv-cot}), multiple-choice formats such as MMLU-STEM (Table \ref{tab:arxiv-cot}), as well as formal mathematics benchmarks like miniF2F (Table \ref{tab:arxiv_atp}).

It should be noted, however, that this conclusion is provisional and subject to certain caveats. Our investigation does not yet address:

\begin{itemize}[topsep=0pt]
    \item How arXiv-sourced data might influence math-related capabilities not covered in the present study, for instance, the informalization of theorems—i.e., translating formal statements or proofs into informal language;
    \item The interaction between arXiv data and other types of training data;
    \item Whether larger model scales could reveal latent benefits of arXiv pre-training.
\end{itemize}

As such, more in-depth analysis is warranted and is reserved for future research.

\subsection{Insights on Reinforcement Learning}
\subsubsection{Moving Towards a Unified Framework} This section seeks to present a consolidated framework for examining a variety of training paradigms, including SFT, RFT, DPO, PPO, GRPO, and extends this with experiments exploring variables within the unified structure. In general, the parameter gradient for any training method can be formalized as:

\begin{equation}
    \nabla_{\theta}\mathcal{J}_{\textcolor{red}{\mathcal{A}}}(\theta) = \mathbb{E}[\underbrace{(q,o) \sim \textcolor{red}{\mathcal{D}}}_{Data \ Source}]\left( \frac{1}{|o|} \sum_{t=1}^{|o|}  \underbrace{GC_{{\mathcal{A}}}(q, o, t, \textcolor{red}{\pi_{{rf}}})}_{Gradient \ Coefficient}  \nabla_{\theta}\log \pi_{\theta}(o_t | q, o_{<t})\right).
\label{eq:objective}
\end{equation}

This framework involves three primary elements: (1) the \textit{Data Source $\mathcal{D}$}, which specifies where training samples are drawn from; (2) the \textit{Reward Function $\pi_{{rf}}$}, responsible for generating the reward feedback used in training; (3) the \textit{Algorithm $\mathcal{A}$}, which integrates the training instances and associated rewards, yielding a gradient coefficient $GC$ that modulates the strength of penalties or rewards applied during learning. We examine a range of leading approaches structured within this cohesive paradigm:

\begin{itemize}
    \item  \textbf{Supervised Fine-tuning (SFT)}: This approach involves adjusting a pretrained model using SFT data curated by humans.
    \item \textbf{Rejection Sampling Fine-tuning (RFT)}: In RFT, the model first undergoes SFT, after which it is further adapted using filtered model outputs generated in response to SFT questions, with filtering based on answer correctness.
    \item \textbf{Direct Preference Optimization (DPO)}: DPO builds upon the SFT model by additional fine-tuning, leveraging extended outputs generated by the SFT model and employing a pairwise DPO loss function.
    \item \textbf{Online Rejection Sampling Fine-tuning (Online RFT)}: Unlike RFT, Online RFT begins by setting the policy model to the SFT model’s parameters and refines it using augmented outputs sampled from the current version of the policy in real-time.
    \item \textbf{PPO/GRPO}: In this method, the SFT model initializes the policy model, which is then reinforced using outputs produced by the policy itself during training.
\end{itemize}

A comprehensive summary of the components within these methodologies can be found in Table \ref{tab:data_policy}. For an in-depth explanation of the derivation steps, please consult Appendix \ref{app:analysis-rl}.

\paragraph{Analysis of Data Source} We categorize the origin of the data into two distinct types: online and offline sampling. Online sampling refers to data collected dynamically from the exploration of the currently-updated training policy, whereas offline sampling utilizes data gathered from the sampling results of the initial SFT model. Both RFT and DPO employ the offline data strategy, while Online RFT and GRPO utilize online data sampling.

From Figure \ref{fig:rl-analysis}, it is evident that Online RFT achieves substantially higher performance than RFT across two benchmark datasets. In particular, Online RFT displays performance on par with RFT during the early phases of training, but later secures a clear lead, underscoring the advantages of using online sampling. This observation aligns with expectations: during initial training iterations, the behaviors of the actor and the SFT model are closely aligned, resulting in only slight variations in the sampled data. As training progresses, however, data derived from the actor increasingly diverges, rendering the benefits of real-time sampling more pronounced.

\paragraph{Analysis of Gradient Coefficient} The process by which these algorithms convert reward signals into gradient coefficients for model updates differs in crucial ways. In our experimental setup, we distinguish the reward functions as either `Rule'-based or `Model'-based. `Rule' designates reward assignment grounded in evaluating the factual correctness of responses, whereas `Model' indicates that responses are scored by a dedicated reward model, itself trained with supervision derived from the rule-based judgments. The formulations provided in Equations \ref{eq:GC-RFT} and \ref{eq:GC-GRPO} showcase a notable distinction: GRPO employs the reward model to modulate its gradient coefficient, thereby enabling nuanced reinforcement or penalization in accordance with the magnitude of the reward. In contrast, Online RFT does not possess this capacity; it simply enforces consistent reinforcement for all correct responses and omits penalties for incorrect ones, regardless of individual differences.

As illustrated in Figure \ref{fig:rl-analysis}, GRPO achieves better performance than online RFT, underscoring the advantage of modifying positive and negative gradient coefficients. Moreover, GRPO+PS outperforms GRPO+OS, suggesting that leveraging step-aware, fine-grained gradient coefficients is beneficial. We also investigate iterative RL, conducting two iterations in our experiments. Figure \ref{fig:iter-rl} reveals that iterative RL yields noticeable performance gains, particularly during the initial iteration.

\subsubsection{Why Does RL Enhance Performance?} This study employs reinforcement learning using a selected portion of the instruction tuning dataset, resulting in notable performance improvements over the instruction-tuned model. To clarify the underlying reasons for RL’s effectiveness, we assess Pass@K and Maj@K accuracy for both Instruct and RL models on two distinct benchmarks. Figure \ref{fig:maj-pass} indicates that while RL boosts Maj@K, Pass@K remains unchanged. These results imply that RL primarily augments the robustness of the output distribution; put differently, \textbf{the observed gains appear to stem from increasing the likelihood of generating correct answers among the TopK candidates, rather than elevating the model’s fundamental problem-solving ability.} This observation aligns with findings from \citep{wang2023making}, who reported a \textbf{misalignment problem} affecting reasoning performance in SFT models, and demonstrated that preference alignment techniques can improve the reasoning abilities of such models \citep{yuan2023rrhf,song2023preference,wang2023making}.

\subsubsection{Approaches for More Effective RL}
\label{sec:effec_rl}
Our results show that RL is highly effective on mathematical reasoning problems. Furthermore, we offer a unified framework for interpreting various representative training strategies, wherein all methods are viewed as either explicit or reduced forms of RL. As formalized in Equation \ref{eq:objective}, three principal elements are involved: Data Source, Algorithm, and Reward Function. We outline prospective research avenues concerning each of these components.

\paragraph{Data Source} The data source underpins every training method. For RL, it specifically refers to unlabeled prompts with outputs generated from the policy model. In this work, we restrict our data to questions from the instruction tuning phase and employ simple nucleus sampling to generate outputs. We hypothesize this limitation partly accounts for our RL pipeline primarily boosting Maj@K scores. Moving forward, we intend to expand our RL pipeline to handle out-of-distribution prompts and integrate \textbf{advanced sampling (decoding) strategies} such as tree-search-based methods \citep{yao2023tree}. Additionally, improvements in \textbf{efficient inference techniques} \citep{xia-etal-2023-speculative,leviathan2023fast,kwon2023efficient,xia2024unlocking} will be pivotal, as they impact how efficiently the policy model explores the solution space.

\paragraph{Algorithms} Algorithms utilize the data and the reward signal to compute the gradient coefficient necessary for updating model parameters. According to Equation \ref{eq:objective}, current methods operate under the assumption that the reward function signal can be fully \textbf{TRUST}ed to either promote or demote the conditional probability associated with specific tokens. Nevertheless, the reliability of the reward signal cannot always be taken for granted, particularly when dealing with highly intricate tasks. For instance, even the PRM800K datasets \citep{lightman2023let}—which have undergone thorough annotation by expert annotators—still contain roughly 20\% annotation errors\footnote{\url{https://github.com/openai/prm800k/issues/12\#issuecomment-1728491852}}. Therefore, it is critical to investigate reinforcement learning algorithms that are inherently resilient to noisy or imperfect reward feedback. We anticipate that alignment approaches following the \textbf{WEAK-TO-STRONG} framework \citep{burns2023weak} may fundamentally transform how learning algorithms operate.

\paragraph{Reward Function} The reward function serves as the foundational training signal in RL, and is typically represented by a neural reward model. We identify three central avenues for the advancement of reward models: 1) \textbf{Strengthening the generalization capabilities of the reward model.} To make meaningful progress, reward models need to generalize effectively to out-of-distribution inputs and sophisticated decoding results; without this, RL could simply reinforce the LLM’s current distribution rather than enhance its core skills; 2) \textbf{Representing the uncertainty within the reward model.} Capturing this uncertainty may bridge the gap between the limitations of weak reward models and the requirements of weak-to-strong learning algorithms; 3) \textbf{Constructing process reward models that deliver high-quality signals efficiently,} particularly those offering detailed feedback throughout the reasoning process \citep{lightman2023let,wang2023math}.

\section{Conclusion, Limitation, and Future Work}
In this paper, we introduce DeepSeekMath, which surpasses all currently available open-source models on the competition-level MATH benchmark and nearly matches the results attained by closed-source models. DeepSeekMath~builds upon the DeepSeek-Coder-v1.5 7B foundation, undergoing continual training for 500B tokens. A substantial portion of this dataset consists of 120B mathematics-related tokens extracted from Common Crawl. Our comprehensive ablation studies reveal that web-based sources present valuable opportunities for curating high-quality mathematical datasets, whereas arXiv appears to be less beneficial than originally anticipated. Furthermore, we present Group Relative Policy Optimization (GRPO)—an adaptation of Proximal Policy Optimization (PPO)—which enhances mathematical reasoning ability while reducing memory usage. Empirical findings confirm that GRPO remains effective, even as DeepSeekMath-Instruct 7B demonstrates strong performance on benchmark tasks. We also propose a unified framework for contextualizing various reinforcement learning approaches and identify several promising avenues for more efficient reinforcement learning.

Despite DeepSeekMath~showing strong results on quantitative reasoning benchmarks, its abilities in the domains of geometry and formal theorem proving remain limited compared to leading closed-source models. For example, during our preliminary evaluation, DeepSeekMath~struggled with questions concerning triangles and ellipses, suggesting possible biases stemming from training data selection during pre-training and fine-tuning stages. Moreover, given the size constraints of the model, DeepSeekMath~is less proficient than GPT-4 in few-shot settings: while GPT-4 demonstrates notable gains with few-shot prompts, DeepSeekMath~performs comparably across both zero-shot and few-shot scenarios. Looking ahead, we plan to refine our data selection strategies to construct a more robust and high-quality pre-training dataset. Additionally, we intend to further investigate the directions outlined in Section \ref{sec:effec_rl} to advance reinforcement learning methods for large language models.